%% ===========================================================================
%%  Caching vs. Offloading vs. Structural Locality:
%%  A Head-to-Head Study of DEX and DART on RDMA Disaggregated Memory
%%
%%  PVLDB / SIGMOD (acmart, sigconf) style.  Compile on Overleaf or with:
%%     pdflatex dex_vs_dart && bibtex dex_vs_dart && pdflatex x2
%%  Figures live in ../compare_plots/ (see \graphicspath below); they are the
%%  PNGs emitted by ../compare_dex_dart.py over the committed sweep CSVs.
%% ===========================================================================
\documentclass[sigconf,nonacm]{acmart}

\usepackage{graphicx}
\usepackage{subcaption}
\usepackage{booktabs}
\usepackage{balance}

% Pull figures straight from the comparison-plot directory so they stay in sync
% with compare_dex_dart.py.  (Copy them into ./figures/ if you move the paper.)
\graphicspath{{../compare_plots/}{C:/Users/abhir/Downloads/REV/compare_plots/}{./figures/}%
  {../dex/build/results/qload/plots/}%
  {C:/Users/abhir/Downloads/REV/dex/build/results/qload/plots/}}

\settopmatter{printacmref=false}
\renewcommand\footnotetextcopyrightpermission[1]{}
\pagestyle{plain}

\begin{document}

\title{REV : Caching, Offloading, or Structural Locality?\\
A Head-to-Head Study of DEX and DART on RDMA Disaggregated Memory}

\author{}
\affiliation{%
  \institution{}
  \country{}
}

\begin{abstract}
Disaggregated memory (DM) splits a server into a compute pool of CPU-rich,
memory-poor nodes and a memory pool of CPU-poor, memory-rich nodes connected by
RDMA. Two recent ordered indexes attack the resulting remote-access bottleneck
from opposite directions. \emph{DART} is a lock-free two-layer hashed Adaptive
Radix Tree that minimizes remote round trips \emph{structurally}---an Express
Skip Table and adaptive hashed nodes keep every operation to a handful of
one-sided reads while the memory node does \emph{zero} CPU work. \emph{DEX} is a
path-aware B$^+$-tree that instead spends two \emph{tunable} resources: a
compute-side directory \emph{cache} and optional RPC \emph{offloading} that runs
index logic on memory-node CPUs to ship results instead of pages. We build a
1-to-1 harness that runs both systems on the same RDMA cluster, working set,
thread budget, and latency instrumentation, and ask a single question:
\emph{when, and by what mechanism, does a tunable caching/offloading index catch
a structurally local one?} Across point lookups and range scans under uniform
and Zipfian-0.99 access, we sweep cache (64--512\,MB), the offload switch, and
the number of memory-node service threads (\texttt{memThreadCount}). We find
that (i)~pure caching alone rarely catches DART---it pays a full page read on
every miss; (ii)~offloading is the mechanism that lets DEX overtake DART,
collapsing per-operation RDMA reads from $2.3$--$22$ to $\approx0$ by shipping an
8\,B value instead of a 0.5--1\,KB page; (iii)~the crossover is workload-shaped:
under Zipfian lookups DEX leads at \emph{every} cache size, while uniform range
scans need the largest cache and still trail in tail latency; and (iv)~offloading
moves a measurable slice of useful work onto memory-node CPUs, where extra
service threads buy throughput only until the network or compute node becomes the
bottleneck. Finally, (v)~we measure remote cost \emph{per operation} (misses and
RDMA reads per op), a fixed-denominator view that makes offloading's effect
directly comparable across cache and explains why its speedup is nearly
independent of the key distribution. We frame these as a portrait of three
operating points and discuss when each design is the right one.
\end{abstract}

\keywords{Disaggregated memory, RDMA, ordered index, caching, compute offloading,
Adaptive Radix Tree, B$^+$-tree}

\maketitle

% ===========================================================================
\section{Introduction}
% ===========================================================================
Disaggregated memory (DM) architectures decouple computing and memory into
separate pools joined by a high-speed interconnect such as RDMA or
CXL~\cite{dart,smart,sherman}. Compute nodes hold powerful CPUs but little
local memory; memory nodes hold large capacity but few or no usable cycles. The
central design tension is well known: a one-sided RDMA operation costs
$\sim$2\,$\mu$s---an order of magnitude more than local memory---and RDMA
bandwidth (100--200\,Gbps) is far below local memory bandwidth. A DM index must
therefore minimize \emph{both} the number of round trips and the bytes moved per
operation, while remaining scalable and crash-consistent under many concurrent
clients.

Two recent systems answer this challenge with fundamentally different
philosophies, and this paper places them side by side.

\paragraph{DART: minimize round trips structurally.}
DART~\cite{dart} is a lock-free two-layer hashed ART. A hash-based
\emph{Express Skip Table} provides express paths that skip the upper levels of
the radix tree, halving traversal depth; an \emph{Adaptive Hashed Layout}
confines each in-node search to one small hash bucket instead of a whole node;
and \emph{Decoupled Metadata Organization} embeds node metadata in the parent
pointer so writes become a single failure-atomic 8\,B compare-and-swap. The
result is an index whose remote cost is low \emph{by construction}: operations
take a handful of one-sided reads, and---crucially---the memory node executes
\emph{no} CPU work at all. Every access is served by the NIC; the memory-node
process blocks on a socket for the duration of the run. DART's performance is
therefore largely \emph{insensitive to compute-side cache}.

\paragraph{DEX: spend cache and memory-node CPU.}
DEX is a path-aware B$^+$-tree for DM. It keeps a compute-side directory
\emph{cache} of inner nodes, and---optionally---\emph{offloads} the leaf access
as an RPC executed on a memory-node service thread, returning the 8\,B value
rather than shipping a 0.5--1\,KB page back to the client. DEX thus exposes two
tunable knobs that DART lacks by design: \emph{how much cache} and \emph{whether
to offload}. Offloading deliberately migrates a slice of useful work onto the
memory node, trading remote CPU cycles for network bytes.

\paragraph{The question.}
These are not just two indexes; they are two \emph{positions} on how to spend the
DM resource budget. DART refuses to spend memory-node CPU and leans on structural
locality. DEX is willing to spend both cache and remote CPU to shrink the bytes
on the wire. The natural question---and the one practitioners actually face when
choosing or provisioning a DM index---is:
\begin{quote}\itshape
When does a tunable caching/offloading index catch a structurally local one, by
what mechanism, and at what cost in memory-node CPU?
\end{quote}

\paragraph{This paper.}
We answer empirically with a controlled 1-to-1 comparison
(Section~\ref{sec:method}). We equalize the working set (50\,M 16\,B entries),
the thread budget (32 compute $+$ 4 memory threads), and the latency artifact
(a shared 500\,ns-bucket histogram yielding p99 on \emph{both} systems), and we
sweep three axes that map exactly onto the design knobs:
\begin{itemize}
  \item \textbf{Cache} (64, 128, 256, 512\,MB) --- the caching axis
        (Section~\ref{sec:cache}).
  \item \textbf{Offload} (off / on) --- the result-shipping axis
        (Section~\ref{sec:offload}).
  \item \textbf{\texttt{memThreadCount}} (2, 4, 6) --- the remote-thread axis
        (Section~\ref{sec:memthreads}).
\end{itemize}
all crossed with operation $\in\{$point lookup, range scan$\}$ and distribution
$\in\{$uniform, Zipfian-0.99$\}$.

\paragraph{Contributions and findings.}
\begin{enumerate}
  \item A \emph{fair-comparison harness} for two DM indexes with opposite
        philosophies, including a per-node CPU sampler and a memory-node
        \emph{active-work} tracker that distinguishes ``cores pegged by
        busy-poll'' from ``useful work done'' (Section~\ref{sec:cycles}).
  \item \textbf{Caching alone rarely catches DART.} DEX with offloading
        \emph{off} pays a full page read on every miss; its per-op RDMA reads
        fall only slowly with cache and it never overtakes DART on uniform
        workloads (Section~\ref{sec:cache}).
  \item \textbf{Offloading is the catch-up mechanism.} Turning offload on
        collapses per-op RDMA reads from $2.3$--$22$ down to $\approx0$, and DEX
        overtakes DART: by $1.7\times$ on uniform lookups and $\approx3\times$
        on Zipfian lookups at 512\,MB (Section~\ref{sec:offload}).
  \item \textbf{The crossover is workload-shaped.} Under Zipfian lookups DEX
        leads at \emph{every} cache size; uniform range scans need the most
        cache and still trail DART in tail latency until $\sim$2\,GB
        (extrapolated) (Sections~\ref{sec:cache}--\ref{sec:offload}).
  \item \textbf{Remote threads have a knee.} With offload on, adding
        memory-node service threads raises throughput only while the memory node
        is the bottleneck; past $\sim$4 threads the gain flattens as the network
        or compute node saturates (Section~\ref{sec:memthreads}).
  \item \textbf{A per-operation remote-cost view.} We measure \emph{misses and
        RDMA reads per operation}---a fixed-denominator metric that compares
        cleanly across cache and the offload switch, quantifies how offloading
        converts page reads into RPCs, and explains why offloading's speedup is
        nearly distribution-independent (Section~\ref{sec:missrate}).
\end{enumerate}

% ===========================================================================
\section{Background and the Two Designs}
\label{sec:bg}
% ===========================================================================

\subsection{Disaggregated memory and RDMA}
A DM compute node accesses the memory pool with one-sided \texttt{RDMA READ}/
\texttt{WRITE} (bypassing the remote CPU) and one-sided atomics
(\texttt{RDMA CAS}/\texttt{FAA}) that update a single 8\,B word. The cost model
has two terms that an index must jointly minimize: \emph{round trips} (each
$\sim$2\,$\mu$s, and serialized by node-to-node dependencies during a traversal)
and \emph{bandwidth} (read amplification, e.g.\ fetching a whole node to find one
pointer). A third axis---\emph{where the cycles run}---is usually ignored because
classic one-sided designs leave the memory-node CPU idle; it becomes central once
a system offloads work (Section~\ref{sec:cycles}).

\subsection{DART: structural locality, idle memory node}
DART minimizes the two cost terms in the structure itself. The Express Skip
Table is a DM-friendly extendible hash table whose directory ($<$1\,MB) is cached
at the compute node, so an express node is reached in a \emph{single} remote read
that skips $2^k$ levels of the tree. Within an inner node, the Adaptive Hashed
Layout hashes the next partial key to one 128\,B bucket, so an in-node search
moves one bucket rather than the whole node. Decoupled Metadata Organization
folds node type/length into the 8\,B compound child pointer, giving lock-free,
failure-atomic writes via a single \texttt{RDMA CAS} and removing the extra round
trip that header metadata would cost. The defining systems property for this
study: DART's memory node performs \emph{no} CPU work---it sets up queue pairs,
signals ready, and blocks on a socket until the run ends. All of DART's useful
work lives on the compute node, and its only ``remote cost'' is network round
trips, not remote cycles.

\subsection{DEX: a path-aware cache plus RPC offloading}
DEX is a B$^+$-tree whose compute nodes hold a \emph{directory cache} of inner
(path) nodes. A lookup walks the cached path and, on reaching the leaf level,
either (a)~\emph{reads the leaf page} over one-sided RDMA (offload \emph{off}),
or (b)~\emph{issues an RPC} to a memory-node service (directory) thread that
performs the leaf access locally and returns just the value (offload \emph{on}).
Mode (a) is a pure caching design directly comparable to DART---no memory-node
CPU---but it pays a full page read on every leaf miss. Mode (b) trades a memory
-node CPU cycle for a large reduction in bytes on the wire: instead of a
0.5--1\,KB page, the network carries an 8\,B value. DEX thus has two knobs DART
structurally lacks: cache size and the offload switch. The number of memory-node
service threads, \texttt{memThreadCount}, bounds how much offloaded work the
memory node can absorb in parallel.

\subsection{Why a busy-poll CPU number is misleading}
Both systems busy-poll the NIC: DEX directory threads run
\texttt{while(true)\{pollWithCQ();process();\}} on pinned cores, and compute
workers spin on completion queues. Raw CPU\% is therefore pinned near 100\%
whether or not real work is happening. To measure ``remote load'' meaningfully we
instrument the \emph{active-work fraction}---time spent inside
\texttt{process\_message} divided by directory-thread wall time---alongside a
per-node process/system CPU sampler. This lets us separate \emph{cores pegged} by
the poll from \emph{useful work done}, which is exactly the distinction that
makes offloading visible (Section~\ref{sec:cycles}).

% ===========================================================================
\section{Experimental Methodology}
\label{sec:method}
% ===========================================================================
We run both systems on a 2-node RDMA cluster (Linux, MLNX\_OFED): one compute
node and one memory node. To make the comparison fair we equalize everything that
is not the design under test.

\paragraph{Working set and matrix.}
Both sweeps load 50\,M keys (16\,B entries) and measure 30--50\,M operations.
\[
\begin{array}{ll}
\text{op}    & \in \{\text{point lookup},\ \text{range scan (100 keys)}\}\\
\text{dist}  & \in \{\text{uniform},\ \text{Zipfian-}0.99\}\\
\text{cache} & \in \{64, 128, 256, 512\}\ \text{MB (compute-side)}\\
\end{array}
\]
DEX names map onto DART's: \textsf{range}$\equiv$\textsf{scan},
\textsf{zipfian}$\equiv$\textsf{zipf99}.

\paragraph{Thread budget.}
We match \textbf{32 compute $+$ 4 memory} threads. DART has no memory-side
service threads (its memory node is idle by construction), so ``4 memory
threads'' is a DEX-only knob that we additionally sweep in
Section~\ref{sec:memthreads}. DART runs 32 compute threads; its per-thread cache
budget is recomputed so the \emph{total} cache matches DEX at each point.

\paragraph{Three system curves.}
Because DART has no offload axis, we run DEX in two modes and carry three curves
on every cache plot: \textbf{DART}, \textbf{DEX offload-off} (pure caching, the
apples-to-apples baseline against DART), and \textbf{DEX offload-on}
(\texttt{MANUAL\_PUSHDOWN}, RPC enabled).

\paragraph{Latency artifact.}
We made the two systems emit an identical latency histogram---ops per 500\,ns
bucket plus per-op p50/p90/p99/max---so latency is compared \textbf{p99-to-p99}
rather than DEX-p99 vs.\ DART-average. All latency and load instrumentation
compiles away when disabled; its overhead is one clock pair per op, negligible
against an RDMA round trip.

\paragraph{Plots.}
All figures are produced by a single script over the committed sweep CSVs, so
every number in the text is reproducible from the same source. Throughput is in
Mops/s (higher is better); latency is p99 in $\mu$s (lower is better).

\paragraph{Miss and traffic instrumentation.}
DEX's path-aware cache carries a per-thread, per-level hit/miss counter that
labels each miss inner- or leaf-node and is reset at the start of each measured
run; dividing the miss counts by the operation count yields \emph{misses per
operation}, and the memory-node counters give the read-vs-RPC split of the remote
traffic. A dedicated queue-load sweep (\texttt{summary\_qload.csv}) repeats the
point-lookup and range-scan matrix with memory-node service threads in $\{2,6\}$;
the per-operation figures in Section~\ref{sec:missrate} are drawn from it. The
counter is exercised in the lookup and scan traversals, which is the entire matrix
here.

% ===========================================================================
\section{The Caching Axis: Does More Cache Catch DART?}
\label{sec:cache}
% ===========================================================================
We first vary cache with offload held in each state.
Figure~\ref{fig:thr} shows throughput; Figure~\ref{fig:p99} shows p99 latency.
The DART curve is essentially \emph{flat}: from 64 to 512\,MB its lookup
throughput holds near 2.8--2.9\,Mops and scan near 1.25--1.30\,Mops, with p99
pinned at $\approx$20.5\,$\mu$s (lookup) and $\approx$40.5\,$\mu$s (scan). This
flatness is the structural-locality story: DART's skip table keeps the hot path
resident with a $<$1\,MB directory, so a larger compute-side cache buys it almost
nothing.

\begin{figure*}[t]
  \centering
  \begin{subfigure}[b]{0.245\textwidth}
    \includegraphics[width=\linewidth]{throughput_vs_cache_lookup_uniform.png}
    \caption{Lookup / uniform}
  \end{subfigure}\hfill
  \begin{subfigure}[b]{0.245\textwidth}
    \includegraphics[width=\linewidth]{throughput_vs_cache_lookup_zipf.png}
    \caption{Lookup / Zipf-0.99}
  \end{subfigure}\hfill
  \begin{subfigure}[b]{0.245\textwidth}
    \includegraphics[width=\linewidth]{throughput_vs_cache_scan_uniform.png}
    \caption{Scan / uniform}
  \end{subfigure}\hfill
  \begin{subfigure}[b]{0.245\textwidth}
    \includegraphics[width=\linewidth]{throughput_vs_cache_scan_zipf.png}
    \caption{Scan / Zipf-0.99}
  \end{subfigure}
  \caption{Throughput vs.\ cache. DART (green) is flat---insensitive to
  compute-side cache. DEX offload-off (blue) rises slowly and, on uniform
  workloads, never catches DART. DEX offload-on (red) climbs steeply and
  overtakes DART once the inner path fits in cache; it leads at \emph{every}
  cache size for Zipfian lookups.}
  \label{fig:thr}
\end{figure*}

\paragraph{Pure caching (offload-off) rarely catches DART.}
DEX offload-off pays a full page read on every leaf miss. Its per-op RDMA reads
fall only slowly with cache---for uniform lookups from 6.79 down to 3.34 over the
$8\times$ cache range---and throughput rises correspondingly slowly: from
1.29 to 2.52\,Mops, \emph{still below} DART's $\approx$2.9 at the largest cache.
On uniform scans it crawls from 0.42 to 0.56\,Mops against DART's 1.29---it is
nowhere close. The lone exception is Zipfian, where caching a small hot set is
genuinely effective: DEX offload-off lookups reach 5.06\,Mops at 512\,MB,
\emph{passing} DART's 2.9 once cache $\geq$256\,MB. Even so, Zipfian \emph{scans}
top out at 0.95\,Mops, below DART's 1.30. Table~\ref{tab:thr} summarizes.

\begin{table}[t]
  \centering
  \caption{Throughput (Mops/s) at the cache endpoints (64\,MB $\to$ 512\,MB).
  DART is flat (single value). Bold marks the best system at 512\,MB.}
  \label{tab:thr}
  \small
  \begin{tabular}{llccc}
    \toprule
    op & dist & DART & DEX off & DEX on \\
    \midrule
    lookup & uniform & 2.9 & $1.29\!\to\!2.52$ & $2.17\!\to\!\mathbf{4.93}$ \\
    lookup & zipf    & 2.9 & $2.04\!\to\!5.06$ & $3.30\!\to\!\mathbf{8.66}$ \\
    scan   & uniform & 1.29 & $0.42\!\to\!0.56$ & $0.90\!\to\!\mathbf{1.46}$ \\
    scan   & zipf    & 1.30 & $0.60\!\to\!0.95$ & $1.29\!\to\!\mathbf{2.37}$ \\
    \bottomrule
  \end{tabular}
\end{table}

\paragraph{Tail latency.}
Figure~\ref{fig:p99} (p99 vs.\ cache) and Figure~\ref{fig:tail} (the focused
tail view with crossover extrapolation) tell the same story for the SLO-relevant
number. DART's p99 is a flat horizontal line; DEX-off's p99 falls with cache but
stays above DART on uniform workloads. The tail view fits DEX's falling p99
against $\log_2(\text{cache})$ and extrapolates to where it would meet DART's
flat tail: for uniform scans that crossover is $\sim$2068\,MB and for Zipfian
scans $\sim$1284\,MB---i.e.\ well beyond the swept range, confirming that for
scans DART's tail is hard to beat with cache alone.

\begin{figure*}[t]
  \centering
  \begin{subfigure}[b]{0.245\textwidth}
    \includegraphics[width=\linewidth]{latency_p99_vs_cache_lookup_uniform.png}
    \caption{Lookup / uniform}
  \end{subfigure}\hfill
  \begin{subfigure}[b]{0.245\textwidth}
    \includegraphics[width=\linewidth]{latency_p99_vs_cache_lookup_zipf.png}
    \caption{Lookup / Zipf-0.99}
  \end{subfigure}\hfill
  \begin{subfigure}[b]{0.245\textwidth}
    \includegraphics[width=\linewidth]{latency_p99_vs_cache_scan_uniform.png}
    \caption{Scan / uniform}
  \end{subfigure}\hfill
  \begin{subfigure}[b]{0.245\textwidth}
    \includegraphics[width=\linewidth]{latency_p99_vs_cache_scan_zipf.png}
    \caption{Scan / Zipf-0.99}
  \end{subfigure}
  \caption{p99 tail latency vs.\ cache (lower is better). DART (green) is flat;
  DEX offload-on (red) drops below DART's tail for lookups within the swept
  range, while for scans it approaches but does not cross it before 512\,MB.}
  \label{fig:p99}
\end{figure*}

% ===========================================================================
\section{The Offloading Axis: Result-Shipping vs. Page-Shipping}
\label{sec:offload}
% ===========================================================================
Caching changes \emph{how often} DEX touches the network; offloading changes
\emph{how much it moves when it does}. This is the mechanism that lets DEX
overtake DART.

\paragraph{RDMA reads collapse.}
With offload on, the dominant leaf miss becomes an RPC that returns an 8\,B value
instead of a page. Per-op RDMA reads collapse accordingly. For uniform lookups
they fall from $2.33$ at 64\,MB to $0.0008$ at 512\,MB---essentially zero remote
reads per operation once the inner path is cached and every leaf access is an
RPC. Zipfian lookups behave the same ($1.30\to0.001$). The work has not
disappeared; it has \emph{moved}: \texttt{rpc/op} sits near 1.0 for lookups (one
leaf RPC per op) and $0.95$--$1.37$ for scans (a scan touches several leaves,
which offload collapses into a small number of RPCs plus one scratch read).

\paragraph{DEX overtakes DART.}
The byte reduction translates directly into throughput (Figure~\ref{fig:thr},
red curve). DEX offload-on overtakes DART:
\begin{itemize}
  \item \textbf{Lookup / uniform}: crosses DART between 64 and 128\,MB; reaches
        $4.93$\,Mops at 512\,MB, $\approx1.7\times$ DART.
  \item \textbf{Lookup / Zipfian}: already above DART at 64\,MB ($3.30$ vs.\
        $2.9$) and leads at \emph{every} cache size, reaching $8.66$\,Mops,
        $\approx3\times$ DART---free hot-set hits compound with result-shipping.
  \item \textbf{Scan / uniform}: the hardest case; it catches DART only near
        512\,MB ($1.46$ vs.\ $1.29$), because each scan offloads several leaf
        reads.
  \item \textbf{Scan / Zipfian}: passes DART by $\sim$128\,MB and reaches
        $2.37$\,Mops, $\approx1.8\times$ DART at 512\,MB.
\end{itemize}

\paragraph{Where DEX catches DART, at a glance.}
Figure~\ref{fig:catch} plots the DEX/DART p99 \emph{ratio} against cache; below
the dashed 1.0 line DEX's tail is faster. The offload-on curve dives under 1.0
for lookups within the swept range (DEX-on lookup p99 reaches $14.0$\,$\mu$s
uniform / $12.5$\,$\mu$s Zipfian at 512\,MB, inside DART's $\approx$20\,$\mu$s
operating range), while offload-off stays above 1.0 on uniform workloads---a
compact statement of ``offloading, not caching, is what closes the tail gap.''

\begin{figure*}[t]
  \centering
  \begin{subfigure}[b]{0.245\textwidth}
    \includegraphics[width=\linewidth]{dex_catches_dart_lookup_uniform.png}
    \caption{Lookup / uniform}
  \end{subfigure}\hfill
  \begin{subfigure}[b]{0.245\textwidth}
    \includegraphics[width=\linewidth]{dex_catches_dart_lookup_zipf.png}
    \caption{Lookup / Zipf-0.99}
  \end{subfigure}\hfill
  \begin{subfigure}[b]{0.245\textwidth}
    \includegraphics[width=\linewidth]{dex_catches_dart_scan_uniform.png}
    \caption{Scan / uniform}
  \end{subfigure}\hfill
  \begin{subfigure}[b]{0.245\textwidth}
    \includegraphics[width=\linewidth]{dex_catches_dart_scan_zipf.png}
    \caption{Scan / Zipf-0.99}
  \end{subfigure}
  \caption{Where DEX catches DART: p99 ratio DEX/DART vs.\ cache. Shaded region
  ($<1.0$) is where DEX's tail is faster than DART's. Offload-on (red) crosses
  into the region for lookups; offload-off (blue) largely does not on uniform
  workloads.}
  \label{fig:catch}
\end{figure*}

\begin{figure*}[t]
  \centering
  \begin{subfigure}[b]{0.245\textwidth}
    \includegraphics[width=\linewidth]{tail_latency_lookup_uniform.png}
    \caption{Lookup / uniform}
  \end{subfigure}\hfill
  \begin{subfigure}[b]{0.245\textwidth}
    \includegraphics[width=\linewidth]{tail_latency_lookup_zipf.png}
    \caption{Lookup / Zipf-0.99}
  \end{subfigure}\hfill
  \begin{subfigure}[b]{0.245\textwidth}
    \includegraphics[width=\linewidth]{tail_latency_scan_uniform.png}
    \caption{Scan / uniform}
  \end{subfigure}\hfill
  \begin{subfigure}[b]{0.245\textwidth}
    \includegraphics[width=\linewidth]{tail_latency_scan_zipf.png}
    \caption{Scan / Zipf-0.99}
  \end{subfigure}
  \caption{Focused tail-latency view. DART's flat p99 (green dashed) is the
  target; DEX offload-on (reds, one curve per \texttt{memThreadCount}) falls
  toward it as cache grows. For scans the dotted line extrapolates the DEX-on
  fit to the cache where it would meet DART ($\sim$2068\,MB uniform,
  $\sim$1284\,MB Zipfian).}
  \label{fig:tail}
\end{figure*}

\paragraph{The ``how'' in one line.}
DEX catches DART by \emph{converting page-shipping into result-shipping}
(offload) while keeping the \emph{inner path resident} (cache $\geq$ the inner
working set). It catches \emph{earlier under skew} (free hot-set hits) and
\emph{needs more cache for scans} (more leaves per op). DART wins precisely when
you cannot or should not spend memory-node CPU: it stays flat-fast with zero
remote load, most visibly on uniform scans.

% ===========================================================================
\section{Per-Operation Remote Cost}
\label{sec:missrate}
% ===========================================================================
Section~\ref{sec:offload} showed \emph{that} offloading collapses DEX's remote
reads; here we quantify the collapse \emph{per operation}. Because the denominator
is the operation count---identical whether offloading is on or off---this view
compares cleanly across cache and the offload switch. DEX's path-aware cache
counts, at each level descended, whether the next node is resolved from the local
cache or fetched from the memory node; dividing the fetch counts by the number of
operations gives \emph{misses per operation}, and the memory-node counters give
the read-vs-RPC split of that remote traffic.

\subsection{Misses per operation across the sweep}
Tables~\ref{tab:mpo-lookup} and~\ref{tab:mpo-range} give misses per operation.
Cache growth lowers every entry; skew lowers it further (Zipf keeps a hot set
resident); offloading roughly halves it for lookups and cuts it $\sim$2.7$\times$
for scans. The inner/leaf split (Figure~\ref{fig:mpostack}) exposes the
mechanism and doubles as a correctness check: a no-offload lookup carries almost
exactly \emph{one} leaf miss per operation (the counter reads
$0.9999$--$0.9968$), which is right---each lookup touches one leaf and the
$\sim$1.9\,GB leaf level never fits---while offloading drives that leaf band to
$\approx0$, because the leaf is served inside the RPC and never fetched as a node.
Scans keep an entry-leaf component but still shrink sharply.

\begin{table}[t]
  \centering
  \caption{Misses per operation, point lookup (2 cores; 6 cores identical to
  within $0.1\%$).}
  \label{tab:mpo-lookup}
  \small
  \begin{tabular}{r rrrr}
    \toprule
    cache (MB) & uni off & uni on & zipf off & zipf on \\
    \midrule
    64  & 6.79 & 3.33 & 4.18 & 2.11 \\
    128 & 5.66 & 2.12 & 3.20 & 1.27 \\
    256 & 4.52 & 1.01 & 2.33 & 0.63 \\
    512 & 3.34 & 1.00 & 1.54 & 0.54 \\
    \bottomrule
  \end{tabular}
\end{table}

\begin{table}[t]
  \centering
  \caption{Misses per operation, range scan (2 cores).}
  \label{tab:mpo-range}
  \small
  \begin{tabular}{r rrrr}
    \toprule
    cache (MB) & uni off & uni on & zipf off & zipf on \\
    \midrule
    64  & 22.35 & 8.18 & 15.30 & 5.33 \\
    128 & 20.94 & 6.98 & 13.43 & 4.23 \\
    256 & 19.22 & 5.68 & 11.46 & 3.20 \\
    512 & 16.88 & 4.26 & \phantom{0}9.28 & 2.22 \\
    \bottomrule
  \end{tabular}
\end{table}

\begin{figure*}[t]
  \centering
  \begin{subfigure}[b]{0.245\textwidth}
    \includegraphics[width=\linewidth]{misses_per_op_vs_cache_lookup_uniform.png}
    \caption{Lookup / uniform}
  \end{subfigure}\hfill
  \begin{subfigure}[b]{0.245\textwidth}
    \includegraphics[width=\linewidth]{misses_per_op_vs_cache_lookup_zipfian.png}
    \caption{Lookup / Zipf-0.99}
  \end{subfigure}\hfill
  \begin{subfigure}[b]{0.245\textwidth}
    \includegraphics[width=\linewidth]{misses_per_op_vs_cache_range_uniform.png}
    \caption{Scan / uniform}
  \end{subfigure}\hfill
  \begin{subfigure}[b]{0.245\textwidth}
    \includegraphics[width=\linewidth]{misses_per_op_vs_cache_range_zipfian.png}
    \caption{Scan / Zipf-0.99}
  \end{subfigure}
  \caption{Total misses per operation vs.\ cache (fixed denominator; 2 cores;
  shared 0--25 scale). Offloading (blue) is $\le$ no-offloading (grey) in every
  panel; both fall with cache and skew.}
  \label{fig:mpototal}
\end{figure*}

\begin{figure*}[t]
  \centering
  \begin{subfigure}[b]{0.245\textwidth}
    \includegraphics[width=\linewidth]{misses_per_op_stacked_lookup_uniform.png}
    \caption{Lookup / uniform}
  \end{subfigure}\hfill
  \begin{subfigure}[b]{0.245\textwidth}
    \includegraphics[width=\linewidth]{misses_per_op_stacked_lookup_zipfian.png}
    \caption{Lookup / Zipf-0.99}
  \end{subfigure}\hfill
  \begin{subfigure}[b]{0.245\textwidth}
    \includegraphics[width=\linewidth]{misses_per_op_stacked_range_uniform.png}
    \caption{Scan / uniform}
  \end{subfigure}\hfill
  \begin{subfigure}[b]{0.245\textwidth}
    \includegraphics[width=\linewidth]{misses_per_op_stacked_range_zipfian.png}
    \caption{Scan / Zipf-0.99}
  \end{subfigure}
  \caption{Misses per operation split into inner (blue) vs.\ leaf (amber); each
  bar-pair is off\,$|$\,on per cache. Offloaded lookups drive the leaf band to
  $\approx0$ (the leaf is served inside the RPC); scans retain an entry-leaf
  component.}
  \label{fig:mpostack}
\end{figure*}

\subsection{Composition of the traffic, and the role of cores}
No-offload remote traffic is 100\% one-sided reads; offloading swaps a slab of
reads for a handful of RPCs (Figure~\ref{fig:rtsplit}). For lookups at
$\ge256$\,MB the read band nearly vanishes and the operation is essentially one
RPC (uniform: $4.5$ reads $\to\approx1.0$ total). Crucially, per-operation traffic
is a \emph{compute-side} quantity: sweeping memory-node cores from 2 to 6 leaves
it unchanged---the curves overlap in Figure~\ref{fig:rtperop}, and misses per op
differ by $<0.1\%$. Cores act on the \emph{rate}, not the per-op demand:
multiplying traffic by throughput, the aggregate remote-op rate delivered to the
memory node \emph{does} separate by core count (Figure~\ref{fig:rtrate}), because
more service threads let throughput climb. This is the queue-load reading of the
knee in Section~\ref{sec:memthreads}.

\begin{figure*}[t]
  \centering
  \begin{subfigure}[b]{0.245\textwidth}
    \includegraphics[width=\linewidth]{remote_traffic_per_op_lookup_uniform.png}
    \caption{Lookup / uniform}
  \end{subfigure}\hfill
  \begin{subfigure}[b]{0.245\textwidth}
    \includegraphics[width=\linewidth]{remote_traffic_per_op_lookup_zipfian.png}
    \caption{Lookup / Zipf-0.99}
  \end{subfigure}\hfill
  \begin{subfigure}[b]{0.245\textwidth}
    \includegraphics[width=\linewidth]{remote_traffic_per_op_range_uniform.png}
    \caption{Scan / uniform}
  \end{subfigure}\hfill
  \begin{subfigure}[b]{0.245\textwidth}
    \includegraphics[width=\linewidth]{remote_traffic_per_op_range_zipfian.png}
    \caption{Scan / Zipf-0.99}
  \end{subfigure}
  \caption{Remote traffic \emph{per operation}, four series (offload off/on
  $\times$ 2/6 cores). The 2- and 6-core markers overlap: per-op traffic is
  compute-side, set by cache misses, not by memory-node threads. Offloading is the
  lever.}
  \label{fig:rtperop}
\end{figure*}

\begin{figure*}[t]
  \centering
  \begin{subfigure}[b]{0.245\textwidth}
    \includegraphics[width=\linewidth]{remote_traffic_split_lookup_uniform.png}
    \caption{Lookup / uniform}
  \end{subfigure}\hfill
  \begin{subfigure}[b]{0.245\textwidth}
    \includegraphics[width=\linewidth]{remote_traffic_split_lookup_zipfian.png}
    \caption{Lookup / Zipf-0.99}
  \end{subfigure}\hfill
  \begin{subfigure}[b]{0.245\textwidth}
    \includegraphics[width=\linewidth]{remote_traffic_split_range_uniform.png}
    \caption{Scan / uniform}
  \end{subfigure}\hfill
  \begin{subfigure}[b]{0.245\textwidth}
    \includegraphics[width=\linewidth]{remote_traffic_split_range_zipfian.png}
    \caption{Scan / Zipf-0.99}
  \end{subfigure}
  \caption{Composition of remote traffic per op: direct RDMA reads (blue) vs.\ RPC
  pushdowns (amber), off\,$|$\,on per cache (2 cores). Offloading converts reads
  into a much smaller number of RPCs.}
  \label{fig:rtsplit}
\end{figure*}

\begin{figure*}[t]
  \centering
  \begin{subfigure}[b]{0.245\textwidth}
    \includegraphics[width=\linewidth]{remote_traffic_rate_lookup_uniform.png}
    \caption{Lookup / uniform}
  \end{subfigure}\hfill
  \begin{subfigure}[b]{0.245\textwidth}
    \includegraphics[width=\linewidth]{remote_traffic_rate_lookup_zipfian.png}
    \caption{Lookup / Zipf-0.99}
  \end{subfigure}\hfill
  \begin{subfigure}[b]{0.245\textwidth}
    \includegraphics[width=\linewidth]{remote_traffic_rate_range_uniform.png}
    \caption{Scan / uniform}
  \end{subfigure}\hfill
  \begin{subfigure}[b]{0.245\textwidth}
    \includegraphics[width=\linewidth]{remote_traffic_rate_range_zipfian.png}
    \caption{Scan / Zipf-0.99}
  \end{subfigure}
  \caption{Aggregate remote-traffic \emph{rate} (throughput $\times$ ops/op). Here
  2 vs.\ 6 cores diverge: more service threads let throughput climb, delivering
  more remote ops/s to the memory node.}
  \label{fig:rtrate}
\end{figure*}

\subsection{Why offloading helps both distributions by the same factor}
The offload/no-offload throughput ratio is nearly identical for uniform and Zipf
(Table~\ref{tab:ratio}) even though absolute throughput differs $\sim$2$\times$.
The misses-per-op data explains why. In this regime throughput is
traffic-bound---for no-offload, throughput$\,\times\,$misses-per-op is nearly
constant ($\approx7.7$ across the lookup sweep). Write misses per op as a product
of a \emph{locality} factor set by the distribution and a \emph{structural} factor
set by offloading. The distribution scales the offload-on and offload-off curves
by the \emph{same} factor, so it cancels in their ratio, leaving only the offload
term:
\[
\frac{\text{thr}_\text{on}}{\text{thr}_\text{off}}
\approx\frac{\text{misses}_\text{off}}{\text{misses}_\text{on}}
=\frac{M_0\,L(\text{dist})\,F_\text{off}}{M_0\,L(\text{dist})\,F_\text{on}}
=\frac{F_\text{off}}{F_\text{on}} .
\]
Empirically the Zipf/uniform misses-per-op ratio at 64\,MB is $0.62$ (off) and
$0.63$ (on)---the same factor that cancels. A residual constant---an RPC costs the
memory node slightly more to serve than a one-sided read---damps the throughput
ratio below the raw traffic ratio, but it is a property of the RPC-vs-RDMA
mechanism, not the distribution, so the matching survives.

\begin{table}[t]
  \centering
  \caption{Throughput ratio (offload on\,$/$\,off), point lookup, 2 cores. Nearly
  identical across distributions despite a $\sim$2$\times$ absolute gap.}
  \label{tab:ratio}
  \small
  \begin{tabular}{r rr}
    \toprule
    cache (MB) & uniform & zipf \\
    \midrule
    64  & 1.66 & 1.59 \\
    128 & 1.85 & 1.74 \\
    256 & 1.74 & 1.73 \\
    512 & 1.48 & 1.45 \\
    \bottomrule
  \end{tabular}
\end{table}

% ===========================================================================
\section{The Remote-Thread Axis: \texttt{memThreadCount}}
\label{sec:memthreads}
% ===========================================================================
Offloading only helps if the memory node can serve the RPCs. Each directory
thread serves RPCs serially while it busy-polls, so the number of service
threads, \texttt{memThreadCount}, bounds offloaded throughput. We sweep it over
$\{2,4,6\}$ with offload off and on across the full matrix.
Figure~\ref{fig:mt} plots throughput vs.\ \texttt{memThreadCount} at the largest
cache; Figure~\ref{fig:mtcache} plots throughput vs.\ cache with one curve per
thread count.

\paragraph{Offload-off is the control---and it is flat.}
With offload off the memory node does no CPU work, so throughput is independent
of \texttt{memThreadCount}: uniform lookups read $1.17$--$1.18$\,Mops at 64\,MB
whether 2, 4, or 6 threads are configured. This is the guardrail: any movement in
the offload-off curve would indicate a confound, and there is none.

\paragraph{Offload-on has a knee.}
With offload on, adding service threads adds parallel memory-node capacity, so
offloaded ops queue less. At 256\,MB, uniform lookups rise from $2.98$\,Mops
(2 threads) to $4.46$ (4) to $4.73$ (6); at 512\,MB from $3.33$ to $4.62$ to
$4.81$. The gain is steep from 2 to 4 threads and then \emph{flattens}: past
$\sim$4 threads the bottleneck has moved off the memory node---onto the network
(NIC message rate) or the compute node (it cannot issue RPCs faster)---so extra
threads buy little. That inflection is the practical answer to ``how many memory
-node threads should I provision for offloading'': about four for this workload,
beyond which the budget is better spent elsewhere.

\begin{figure}[t]
  \centering
  \begin{subfigure}[b]{0.235\textwidth}
    \includegraphics[width=\linewidth]{memthreads_lookup_uniform.png}
    \caption{Lookup / uniform}
  \end{subfigure}\hfill
  \begin{subfigure}[b]{0.235\textwidth}
    \includegraphics[width=\linewidth]{memthreads_lookup_zipf.png}
    \caption{Lookup / Zipf-0.99}
  \end{subfigure}
  \caption{Throughput vs.\ \texttt{memThreadCount} (offload-on, largest cache).
  Throughput rises while the memory node is the bottleneck, then flattens past
  $\sim$4 threads.}
  \label{fig:mt}
\end{figure}

\begin{figure*}[t]
  \centering
  \begin{subfigure}[b]{0.245\textwidth}
    \includegraphics[width=\linewidth]{memthreads_cache_lookup_uniform.png}
    \caption{Lookup / uniform}
  \end{subfigure}\hfill
  \begin{subfigure}[b]{0.245\textwidth}
    \includegraphics[width=\linewidth]{memthreads_cache_lookup_zipf.png}
    \caption{Lookup / Zipf-0.99}
  \end{subfigure}\hfill
  \begin{subfigure}[b]{0.245\textwidth}
    \includegraphics[width=\linewidth]{memthreads_cache_scan_uniform.png}
    \caption{Scan / uniform}
  \end{subfigure}\hfill
  \begin{subfigure}[b]{0.245\textwidth}
    \includegraphics[width=\linewidth]{memthreads_cache_scan_zipf.png}
    \caption{Scan / Zipf-0.99}
  \end{subfigure}
  \caption{Joint \texttt{memThreadCount}$\times$cache effect on DEX offload-on
  throughput, with DEX offload-off and DART's flat baseline for reference. More
  threads lift the whole curve while the memory node is the bottleneck; the lift
  shrinks once cache makes the compute/network side the limit.}
  \label{fig:mtcache}
\end{figure*}

\paragraph{Interaction with cache.}
Figure~\ref{fig:mtcache} shows the two knobs are not independent. At small cache
the memory node is busiest (many leaf RPCs), so more threads help most; at large
cache the inner path is resident and the limit shifts to the compute/network
side, so the thread curves bunch together. The two DEX knobs are complementary:
cache decides how \emph{often} you offload, threads decide how fast the memory
node \emph{absorbs} what you offload.

% ===========================================================================
\section{Where the Cycles Live: Remote vs. Compute Load}
\label{sec:cycles}
% ===========================================================================
The three operating points differ not only in throughput but in \emph{where the
useful work runs}---the axis disaggregation is supposed to care about.

\begin{table}[t]
  \centering
  \caption{Where the cycles live. ``Useful'' is the active-work fraction; both
  systems busy-poll, so gross CPU\% is pinned regardless.}
  \label{tab:cycles}
  \small
  \begin{tabular}{lll}
    \toprule
    node / mode & gross CPU\% & useful work \\
    \midrule
    DART memory          & $\approx$0 (blocked) & 0 (passive RDMA target) \\
    DART compute         & high (workers+spin)  & all index logic \\
    DEX memory, off      & $\approx$100\%/thread (spin) & $\approx$0 active \\
    DEX memory, on       & $\approx$100\%/thread (spin) & $>$0, rises with rpc \\
    DEX compute          & high (workers+spin)  & index + RPC issue \\
    \bottomrule
  \end{tabular}
\end{table}

\paragraph{DART and DEX-caching: memory node idle.}
DART's memory node does zero useful work by construction, and DEX offload-off is
the same---caching is one-sided, so the memory-node CPU is idle (busy-polling but
doing nothing). All useful work is on the compute node. This is ``pure''
disaggregation: the memory pool is dumb memory.

\paragraph{DEX-offload: a measured slice migrates.}
Offload deliberately moves useful work onto the memory node. Even though gross
CPU\% was already pinned by the poll, the \emph{active-work fraction} rises above
zero and grows with the RPC rate---this is the only regime where memory-node CPU
load, and hence \texttt{memThreadCount}, matters at all. The
\texttt{rpc/op}$\approx$1 (lookups) to $1.37$ (uniform scans) figures from
Section~\ref{sec:offload} quantify exactly how much work migrates per operation.
Table~\ref{tab:cycles} summarizes the four regimes. The takeaway: DEX buys its
throughput lead by spending a resource---memory-node cycles---that DART refuses
to touch.

% ===========================================================================
\section{Synthesis: Three Operating Points}
\label{sec:synth}
% ===========================================================================
Pulling the axes together, each system occupies a distinct position in a
five-dimensional trade-off space: \emph{Throughput}, \emph{Tail speed}
($1/$p99), \emph{Network thrift} ($1/$RDMA-reads-per-op), \emph{Cache
robustness} (flatness across the cache sweep), and \emph{Disaggregation purity}
(how little memory-node CPU it uses---maximal for DART and DEX-off, reduced for
DEX-on).

\begin{itemize}
  \item \textbf{DART --- the spiky specialist.} Strong on \emph{network thrift},
        \emph{cache robustness}, and \emph{disaggregation purity}; respectable on
        \emph{tail speed}; weak only where it structurally cannot compete---it
        can never trade memory-node CPU for network bytes. Its advantages stand
        out most under Zipfian-small-cache.
  \item \textbf{DEX offload-off --- the same shape, smaller.} It shares DART's
        ``no memory-node CPU'' profile but with worse network thrift (full page
        reads), so it sits inside DART on most axes.
  \item \textbf{DEX offload-on --- the balanced, larger envelope.} It activates
        the latent offload capability, pushes network thrift to the maximum
        (reads/op $\to0$), and takes the throughput/tail lead at adequate
        cache---enclosing the other two under uniform/large-cache---paying for it
        only in disaggregation purity.
\end{itemize}

\paragraph{A decision rule.}
Choose DART when memory-node CPU is unavailable or must stay idle (strict
disaggregation, shared memory pools, untrusted/limited memory nodes), when cache
is scarce, or for scan-heavy uniform workloads where its flat tail is hard to
beat. Choose DEX with offloading when memory-node cycles are spendable and the
inner path fits in cache---especially under skew, where it leads everywhere, or
when network bandwidth is the binding constraint and shipping 8\,B beats shipping
a page.

% ===========================================================================
\section{Discussion and Threats to Validity}
\label{sec:disc}
% ===========================================================================
\paragraph{Absolute throughput and thread count.}
Throughput numbers are absolute and depend on the matched 32/4 thread budget;
the \emph{shapes} (flat DART, slow DEX-off, steep DEX-on, the knee in
\texttt{memThreadCount}) are the robust findings. DART's cache flatness and DEX's
RDMA-read collapse are mechanism-level effects that do not depend on the exact
thread count.

\paragraph{Crossover precision.}
The scan tail crossovers ($\sim$2068\,MB uniform, $\sim$1284\,MB Zipfian) are
log-linear extrapolations beyond the swept range and should be read as ``DART's
scan tail is not caught within practical cache,'' not as exact cache budgets.

\paragraph{Remote-load magnitudes.}
We report the \emph{direction} and the per-op RPC counts that drive memory-node
load; absolute active-work percentages depend on the reporter window and are used
qualitatively (Table~\ref{tab:cycles}).

\paragraph{Single memory node.}
The topology is one compute / one memory node. Multi-node scaling, failure
recovery, and write-heavy/mixed workloads are natural extensions; both systems
provide the crash-consistency primitives (DART's failure-atomic CAS, DEX's RPC
path) to support them.

% ===========================================================================
\section{Related Work}
\label{sec:rel}
% ===========================================================================
DM indexing has produced both tree and hash designs. SMART~\cite{smart} is the
state-of-the-art DM-based ART and caches inner nodes at the compute node, but is
limited by compute-side cache capacity and a FIFO cache policy, and reads whole
nodes (read amplification). Sherman~\cite{sherman} is a write-optimized DM
B$^+$-tree using lock-based concurrency with hierarchical locks. DART~\cite{dart}
removes both the round-trip and bandwidth costs structurally and adds lock-free,
failure-atomic writes. RACE hashing~\cite{race} is the DM-friendly extendible
hash that DART's Express Skip Table builds on. DEX represents the
caching-plus-offloading lineage, where compute pushes work into the memory pool
to save network bytes. Our contribution is not a new index but a \emph{controlled
head-to-head} that isolates the mechanisms---caching, offloading, remote
threads---and shows exactly which one closes the gap and at what cost.

% ===========================================================================
\section{Conclusion}
\label{sec:conc}
% ===========================================================================
DART and DEX embody opposite answers to the disaggregated-memory index question.
DART minimizes remote cost \emph{structurally} and keeps the memory node idle;
DEX spends \emph{cache} and \emph{memory-node CPU} to do the same. Our 1-to-1
study shows that caching alone rarely catches DART, that \emph{offloading} is the
mechanism that lets DEX overtake it---collapsing per-op RDMA reads to near zero
and leading by up to $\approx3\times$ under skew---and that the crossover is
workload-shaped, with uniform scans the hardest case. The advantage is bought
with memory-node cycles whose absorption is gated by a small number of service
threads. The practical lesson is a portrait of three operating points rather than
a single winner: structural locality when the memory node must stay dumb,
caching-plus-offloading when its cycles are yours to spend.

\balance
\begin{thebibliography}{9}
\bibitem{dart} B.~Zhang et al. DART: A Lock-free Two-layer Hashed ART Index for
  Disaggregated Memory. \emph{Proc. ACM Manag. Data} 4(1) (SIGMOD), 2026.
\bibitem{smart} X.~Luo et al. SMART: A High-Performance Adaptive Radix Tree for
  Disaggregated Memory. \emph{OSDI}, 2023.
\bibitem{sherman} Q.~Wang et al. Sherman: A Write-Optimized Distributed B+Tree
  Index on Disaggregated Memory. \emph{SIGMOD}, 2022.
\bibitem{race} P.~Zuo et al. One-sided RDMA-Conscious Extendible Hashing for
  Disaggregated Memory (RACE). \emph{USENIX ATC}, 2021.
\end{thebibliography}

\end{document}
